% =====================================================================
%  config/estilo.tex
%  Capa, folha de rosto, metadados e aparência geral.
%  A equipe não precisa mexer aqui.
% =====================================================================

% --- Lista de integrantes em três formatos ----------------------------
\newcommand{\integrante}[3]{}

\newcommand{\imprimirIntegrantesCapa}{%
  \begingroup
  \renewcommand{\integrante}[3]{##1\par}%
  \listaIntegrantes
  \endgroup}

\newcommand{\imprimirIntegrantesLinha}{%
  \begingroup
  \renewcommand{\integrante}[3]{##1; }%
  \listaIntegrantes
  \endgroup}

\newcommand{\imprimirTabelaIntegrantes}{%
  \begingroup
  \renewcommand{\integrante}[3]{##1 & ##2 & ##3 \\}%
  \noindent
  \begin{tabular}{@{}L{7.4cm}L{3.0cm}L{4.2cm}@{}}
    \toprule
    \textbf{Integrante} & \textbf{Matrícula} & \textbf{Papel no ciclo} \\
    \midrule
    \listaIntegrantes
    \bottomrule
  \end{tabular}%
  \endgroup}

% --- Metadados do abnTeX2 ---------------------------------------------
\titulo{\tituloProjeto}
\autor{\nomeEquipe}
\local{\localEntrega, Piauí}
\data{\dataEntrega}
\instituicao{\instituicaoCurso\par\unidadeCurso\par\programaCurso}
\tipotrabalho{Relatório técnico}
\orientador[Docente:]{\docenteDisciplina}
\preambulo{%
  Relatório técnico do \nomeciclo, apresentado como instrumento da
  \rotuloAP\ da disciplina \nomeDisciplina, \periodoLetivo,
  do \programaCurso\ da \instituicaoCurso.}

% --- Links e metadados do PDF -----------------------------------------
\definecolor{corLink}{HTML}{0B3D91}
\makeatletter
\hypersetup{
  pdftitle    = {\@title},
  pdfauthor   = {\@author},
  pdfsubject  = {\nomeDisciplina},
  pdfkeywords = {IoT}{AIoT}{ESP32}{sistemas embarcados}{UFPI},
  colorlinks  = true,
  linkcolor   = corLink,
  citecolor   = corLink,
  urlcolor    = corLink,
  filecolor   = corLink,
  bookmarksdepth = 4,
  breaklinks  = true,
}
\makeatother

% --- Espaçamentos e listas --------------------------------------------
\setlrmarginsandblock{3cm}{2cm}{*}
\setulmarginsandblock{3cm}{2cm}{*}
\checkandfixthelayout
\setlength{\parindent}{1.3cm}
\setlength{\parskip}{0.2cm}
\setlist{itemsep=1pt, topsep=4pt, parsep=0pt}
\SingleSpacing

% --- Capa customizada --------------------------------------------------
\renewcommand{\imprimircapa}{%
  \begin{capa}
    \centering
    \IfFileExists{figuras/logo-ufpi.png}{%
      \includegraphics[height=2.3cm]{figuras/logo-ufpi.png}\par\vspace{0.6cm}}{}
    {\sffamily
      {\large\MakeUppercase{\instituicaoCurso}\par}
      {\normalsize\MakeUppercase{\unidadeCurso}\par}
      {\normalsize\MakeUppercase{\programaCurso}\par}}
    \vspace{1.6cm}
    {\sffamily\small\MakeUppercase{Equipe \nomeEquipe}\par}
    \vspace{0.5cm}
    {\normalsize\imprimirIntegrantesCapa}
    \vfill
    {\sffamily\bfseries\LARGE\tituloProjeto\par}
    \vspace{0.4cm}
    {\sffamily\large\itshape\subtituloProjeto\par}
    \vspace{1.2cm}
    \rule{0.55\textwidth}{0.6pt}\par
    \vspace{0.6cm}
    {\sffamily\normalsize Relatório técnico da \rotuloAP\par}
    {\sffamily\small\nomeciclo\par}
    \vspace{0.3cm}
    {\small\nomeDisciplina\par}
    {\small\periodoLetivo\ \textbullet\ \docenteDisciplina\par}
    \vfill
    {\large\localEntrega\par}
    {\large\dataEntrega\par}
    \vspace{0.8cm}
  \end{capa}}
